\section{From a counterexample to a strict subset inequality}\label{sec:reduction}

\begin{contract}
\textbf{Assumes:} that \cref{thm:main} has a counterexample. \quad
\textbf{Produces:} an oriented graph $G$, all of whose indegrees are positive, satisfying the
strict inequality
\[
 \abs{U}+\abs{F_G^2(U)}<2\abs{F_G(U)}
 \qquad\text{for every nonempty proper }U\subseteq V(G).
\]
\textbf{Why:} this is the graph that \cref{sec:incompatibility} destroys. The section is the
\emph{first engine}: a minimal counterexample is forced to satisfy a strictly stronger, purely
set-theoretic inequality, and a blowup amplifies the integer margin so that later rounding cannot
lose it.
\end{contract}

Throughout, $F=F_D$ is the \emph{union of out-neighborhoods} in a graph $D$:
\begin{equation}\label{eq:Fdef}
 F(S)=\{y\in V(D):\ \text{there is }x\in S\text{ with }x\to y\},
 \qquad F^k(S)=\underbrace{F(F(\cdots F}_{k}(S)\cdots)).
\end{equation}
Two warnings about \eqref{eq:Fdef}, both used constantly. First, $F(S)$ does \emph{not} subtract
$S$: an arc inside $S$ keeps its head in $F(S)$. Second, $F^2(\{v\})$ does \emph{not} exclude $v$;
what excludes $v$ is the orientation. In an oriented graph there is no two-step walk from $v$ back
to $v$ (that would need $v\to u\to v$, two opposite arcs), so
\begin{equation}\label{eq:N2equals}
 N_2^+(v)=F^2(\{v\})\setminus F(\{v\}),
\end{equation}
with no need to remove $v$ separately. The two operators $F$ and $N_2^+$ are the same object seen
from two angles: $F$ acts on sets of vertices, $N_2^+$ is the special case of a single vertex, with
the first neighborhood subtracted.

The goal of the section is the following.

\begin{proposition}[Counterexample reduction]\label{prop:reduction}
If \cref{thm:main} has a counterexample, then there is a finite nonempty oriented graph $G$ in which
every vertex has positive indegree and
\begin{equation}\label{eq:strict-growth}
 \abs{U}+\abs{F_G^2(U)}<2\abs{F_G(U)}
 \qquad\text{for every }\varnothing\ne U\subsetneq V(G).
\end{equation}
\end{proposition}

\inwords{If the conjecture fails somewhere, then somewhere else there is a graph in which every
nonempty proper set of vertices has its second image strictly smaller than the average of the set
and its first image.}

\begin{planbox}
\begin{enumerate}[leftmargin=1.3em,itemsep=1pt,topsep=1pt]
\item \textbf{Subset deficit} (\cref{lem:deficit}): in a counterexample chosen minimal first by
  order and then by number of arcs, every nonempty proper $S$ satisfies the \emph{strict} deficit
  $\abs{F^2(S)\setminus F(S)}<\abs{F(S)\setminus S}$. The right-hand side is the part of the first
  image that lies outside $S$; that refinement is what makes the rest work.
\item \textbf{Blowup} (\cref{lem:blowup}): replacing every vertex by a transitive tournament on
  $m>\abs{V(D)}$ vertices multiplies the deficit by $m$ and costs at most $n$ in rounding, turning
  an integer deficit (at least $1$) into \eqref{eq:strict-growth}.
\end{enumerate}
\end{planbox}

\begin{proof}[Proof of \cref{prop:reduction}, assuming \cref{lem:deficit} and \cref{lem:blowup}]
Choose a counterexample $D$ of minimum order, and among those of minimum number of arcs.
\Cref{lem:deficit} supplies positive indegrees for $D$ and the deficit
$\abs{F^2(S)\setminus F(S)}<\abs{F(S)\setminus S}$ for every nonempty proper $S$. Apply
\cref{lem:blowup} with any integer $m>\abs{V(D)}$; the resulting $G$ is as required.
\end{proof}

\subsection{The subset deficit}\label{sec:deficit}

\begin{ideabox}
The conjecture says $\abs{N_2^+(v)}<\abs{N_1^+(v)}$ at every vertex of a counterexample. We want a
statement about \emph{sets} of vertices rather than single vertices, and we want it to be strict
with a robust right-hand side. The device that produces it is: delete the arcs from $S$ to a
carefully chosen part of $E=F(S)\setminus S$, and see what minimality says about the vertex that
must now satisfy the inequality. Deleting arcs can only shrink first neighborhoods and can only
shrink second neighborhoods; so the vertex that gains from the deletion must have \emph{lost} some
out-neighbors, and the bookkeeping of that trade is the content of the lemma.
\end{ideabox}

\begin{lemma}[Subset deficit]\label{lem:deficit}
Let $D$ be a finite nonempty oriented graph with $\abs{N_1^+(v)}>\abs{N_2^+(v)}$ for every
$v\in V(D)$, chosen with minimum order and, subject to that, with minimum number of arcs. Then $D$
is strongly connected, every vertex has positive indegree, and
\begin{equation}\label{eq:deficit}
 \abs{F^2(S)\setminus F(S)}<\abs{F(S)\setminus S}
 \qquad\text{for every }\varnothing\ne S\subsetneq V(D).
\end{equation}
\end{lemma}

\inwords{In a smallest counterexample, every proper nonempty set $S$ loses, at its second image,
strictly less than the size of the part of its first image that lies outside $S$.}

\usedin{This is the only place where minimality of the counterexample is used; it feeds
\cref{lem:blowup} and hence \cref{prop:reduction}}

\begin{planbox}
\begin{enumerate}[leftmargin=1.3em,itemsep=1pt,topsep=1pt]
\item \textbf{Clean up the graph} (Step~1): no sink, at least two vertices, strongly
  connected, all indegrees positive.
\item \textbf{Fix $S$ and split the first image} (Step~2): $I=F(S)\cap S$ is the
  \emph{internal} part, $E=F(S)\setminus S$ the \emph{external} part; $E\ne\varnothing$. Define
  $g(Y)=\abs{F(Y)\setminus(I\cup Y)}$ for $Y\subseteq E$, so that the target inequality reads
  $g(E)<\abs E$.
\item \textbf{Build a deletion machine} (Step~3): for $Y\subsetneq E$, delete all arcs
  from $S$ to $E\setminus Y$. Minimality of the arc count gives a vertex $v$ satisfying the
  inequality in the smaller graph; show $v\in S$, and let $Q_v$ be the out-neighbors it loses.
\item \textbf{Do the bookkeeping} (Step~4): the second out-neighborhood of $v$ loses
  exactly a set $\Lambda$ and gains only points of $Q_v\cap F(Y)$; hence the increase of $g$ is
  less than $\abs{Q_v}$.
\item \textbf{Iterate} (Step~5) from $Y=\varnothing$ to $Y=E$, and sum.
\end{enumerate}
\end{planbox}

\begin{proof}
\steplabel{Step 1: the graph is strongly connected.}
A sink satisfies $\abs{N_1^+(v)}=0\le\abs{N_2^+(v)}$, so $D$ has no sink; in particular $D$ is not
a single vertex, since a single vertex is a sink. If $D$ were not strongly connected, its
condensation would have a sink component $C$: a strongly connected component with no arc to
another component. Every out-neighbor, and therefore every two-step out-neighbor, of a vertex of
$C$ lies in $C$. Hence $D[C]$ would be a counterexample of smaller order, contradicting minimality.
So $D$ is strongly connected, and since it has more than one vertex, every vertex has an
in-neighbor.

\steplabel{Step 2: fix a subset and split its first image.}
Let $\varnothing\ne S\subsetneq V(D)$ and put
\begin{equation}\label{eq:IE}
 I=F(S)\cap S,\qquad E=F(S)\setminus S,\qquad F(S)=I\sqcup E,
\end{equation}
and, for $Y\subseteq E$,
\begin{equation}\label{eq:gdef}
 g(Y)=\abs{F(Y)\setminus(I\cup Y)} .
\end{equation}
We have $E\ne\varnothing$: otherwise $S$ would have no arc to $V(D)\setminus S$, contradicting
strong connectivity as $S$ is nonempty and proper. The reader should keep \cref{fig:deficit} in
mind: $I$ is the part of the first image that falls back inside $S$, $E$ the part that falls
outside.

\begin{figure}[t]
\centering
\begin{tikzpicture}[
  svert/.style={circle,draw=blue!50!black,fill=blue!10,inner sep=0pt,minimum size=5mm,font=\scriptsize},
  ivert/.style={circle,draw=blue!50!black,fill=blue!30,inner sep=0pt,minimum size=5mm,font=\scriptsize},
  evert/.style={circle,draw=orange!60!black,fill=orange!28,inner sep=0pt,minimum size=5mm,font=\scriptsize},
  gvert/.style={circle,draw=green!40!black,fill=green!25,inner sep=0pt,minimum size=5mm,font=\scriptsize},
  arc/.style={-{Stealth[length=1.5mm]},draw=black!70},
  lab/.style={font=\scriptsize,text=black!75}
]
% ambient box
\draw[draw=black!45,rounded corners=6pt,fill=black!2] (0,0) rectangle (12.6,5.7);
\node[lab,anchor=north west] at (0.18,5.62) {$V(D)$};
% S
\draw[blue!45!black,fill=blue!7,thick] (3.6,3.0) ellipse (2.9 and 2.25);
\node[lab,text=blue!45!black,anchor=west] at (0.12,3.0) {$S$};
% vertices
\node[svert] (s1) at (1.4,4.1) {};
\node[svert] (s2) at (1.9,1.6) {};
\node[svert] (s3) at (2.6,2.35) {};
\node[ivert] (i1) at (4.9,4.1) {};
\node[ivert] (i2) at (5.35,2.75) {};
\node[evert] (e1) at (7.5,4.6) {};
\node[evert] (e2) at (7.9,1.6) {};
\node[gvert] (g1) at (10.0,4.6) {};
\node[gvert] (g2) at (10.3,1.6) {};
% arcs
\draw[arc] (s1) -- (i1);
\draw[arc] (s3) -- (i2);
\draw[arc] (i2) -- (e1);
\draw[arc] (i1) -- (e1);
\draw[arc] (e1) -- (g1);
\draw[arc] (s2) -- (e2);
\draw[arc] (e2) -- (g2);
% region labels, all clear of nodes and arcs
\node[lab,text=blue!45!black,anchor=south] at (3.1,4.25) {$I=F(S)\cap S$};
\node[lab,text=orange!55!black,anchor=south] at (7.7,4.95) {$E=F(S)\setminus S$};
\node[lab,text=green!35!black,anchor=south] at (10.15,4.95) {$F^2(S)\setminus F(S)$};
% the inequality, on a plate below the ellipse
\node[draw=black!25,fill=white,rounded corners=2pt,inner sep=4pt,font=\scriptsize]
  at (9.1,0.42) {$\abs{F^2(S)\setminus F(S)} \;<\; \abs{F(S)\setminus S}=\abs E$};
\end{tikzpicture}

\caption{The split \eqref{eq:IE} of the first image of $S$. \emph{Look at} the three regions: $I$
 (the first image that stays inside $S$), $E$ (the first image outside $S$), and the region
 $F^2(S)\setminus F(S)$ (second image not already seen in one step). \emph{Why it matters:}
 \cref{lem:deficit} compares the size of the third region with the size of $E$, not with the size
 of $F(S)=I\sqcup E$. The part $I$ is free: it is already ``paid for'' by $S$.}
\label{fig:deficit}
\end{figure}

Note the two extremes: $g(\varnothing)=0$, and $g(E)=\abs{F(E)\setminus(I\cup E)}$. Since
$F(S)=I\sqcup E$ and $F(I)\subseteq F(S)=I\sqcup E$, we shall be able to identify the target
quantity at the end:
\begin{equation}\label{eq:deficit-via-g}
 F^2(S)\setminus F(S)
 =\bigl(F(I)\cup F(E)\bigr)\setminus(I\cup E)
 =F(E)\setminus(I\cup E),
\end{equation}
so that $\abs{F^2(S)\setminus F(S)}=g(E)$.

\steplabel{Step 3: the deletion machine.}
Fix $Y\subsetneq E$. Let $D_Y$ be the graph obtained from $D$ by deleting \emph{all} arcs from $S$
to $E\setminus Y$; all other arcs are kept. Every vertex of $E\setminus Y$ has an in-neighbor in
$S$, because $E\setminus Y\subseteq E=F(S)\setminus S$; so at least one arc is deleted and $D_Y$
has strictly fewer arcs than $D$.

Minimality of the number of arcs (the order is unchanged) says that $D_Y$ is \emph{not} a
counterexample. So there is a vertex $v$ with
\begin{equation}\label{eq:DY-inequality}
 \abs{N_{1,D_Y}^+(v)}\le\abs{N_{2,D_Y}^+(v)}.
\end{equation}
If $v$ had lost no out-neighbor, then $N_{1,D_Y}^+(v)=N_1^+(v)$ while $N_{2,D_Y}^+(v)\subseteq
N_2^+(v)$ (deleting arcs creates no paths), so
$\abs{N_{2,D_Y}^+(v)}\le\abs{N_2^+(v)}<\abs{N_1^+(v)}=\abs{N_{1,D_Y}^+(v)}$, contradicting
\eqref{eq:DY-inequality}. Hence $v$ lost an out-neighbor, which means $v\in S$: only arcs with tail
in $S$ were deleted. Put
\begin{equation}\label{eq:Qv}
 Q_v=F(\{v\})\cap(E\setminus Y)=\{\text{out-neighbors of }v\text{ that }D_Y\text{ does not have}\}.
\end{equation}
Then $Q_v\ne\varnothing$, and $Q_v$ is disjoint from $I\cup Y$ because $Q_v\subseteq E\setminus Y$
and $Y\subseteq E$ is disjoint from $S\supseteq I$. Also $F(\{v\})\subseteq I\cup Y\cup Q_v$,
since $F(\{v\})\subseteq F(S)=I\sqcup E$.

\steplabel{Step 4: what $v$ loses and what it gains.}
\Cref{fig:deletion} shows the four sets involved. We record four facts.

\emph{(4a) The second image after deletion is confined.} Since $v\in S$, its out-neighbors in
$D_Y$ lie in $I\cup Y$; and for such an out-neighbor $x$,
\begin{itemize}[leftmargin=1.4em,itemsep=1pt,topsep=1pt]
\item if $x\in I\subseteq S$, then all arcs out of $x$ into $E\setminus Y$ were deleted, so
  $F_{D_Y}(x)\subseteq F(x)\setminus(E\setminus Y)\subseteq (I\sqcup E)\setminus(E\setminus
  Y)=I\cup Y$;
\item if $x\in Y$, then $F_{D_Y}(x)\subseteq F(x)\subseteq F(Y)$.
\end{itemize}
Hence
\begin{equation}\label{eq:confined}
 F_{D_Y}^2(\{v\})\subseteq I\cup Y\cup F(Y).
\end{equation}

\emph{(4b) New second out-neighbors are first out-neighbors.} A vertex counted in
$N_{2,D_Y}^+(v)$ but not in $N_2^+(v)$ is the end of a two-step path of $D_Y$, hence of $D$; it is
not a first out-neighbor in $D_Y$. So it was a first out-neighbor in $D$, i.e.\ it lies in $Q_v$;
and by \eqref{eq:confined} it lies in $I\cup Y\cup F(Y)$. As $Q_v$ misses $I\cup Y$,
\begin{equation}\label{eq:new}
 N_{2,D_Y}^+(v)\setminus N_2^+(v)\subseteq Q_v\cap F(Y).
\end{equation}

\emph{(4c) The set $\Lambda$ is lost.} Define
\begin{equation}\label{eq:Lambda}
 \Lambda=F(Q_v)\setminus\bigl(I\cup Y\cup Q_v\cup F(Y)\bigr).
\end{equation}
Each $y\in\Lambda$ has $v\to q\to y$ with $q\in Q_v\subseteq F(\{v\})$, so $y\in F^2(\{v\})$. Also
$y\notin F(\{v\})$, because $F(\{v\})\subseteq I\cup Y\cup Q_v$; and $y\ne v$, because $y\in F(Q_v)$
and an arc $q\to v$ would be opposite to the arc $v\to q$. Hence
\begin{equation}\label{eq:Lambda-sub}
 \Lambda\subseteq N_2^+(v).
\end{equation}
On the other hand $\Lambda\cap(I\cup Y\cup F(Y))=\varnothing$ by definition, so \eqref{eq:confined}
gives $\Lambda\cap N_{2,D_Y}^+(v)=\varnothing$: every point of $\Lambda$ was counted before the
deletion and is not counted after it.

\emph{(4d) The change of $g$ is exact.} Since $F(Y\cup Q_v)=F(Y)\cup F(Q_v)$,
\begin{align}
 F(Y\cup Q_v)\setminus(I\cup Y\cup Q_v)
 &=\bigl(F(Y)\setminus(I\cup Y\cup Q_v)\bigr)
   \;\sqcup\;\bigl(F(Q_v)\setminus(I\cup Y\cup Q_v\cup F(Y))\bigr)\notag\\
 &=\bigl(F(Y)\setminus(I\cup Y)\bigr)\setminus Q_v\;\;\sqcup\;\;\Lambda,
 \label{eq:g-change}
\end{align}
using $Q_v\cap(I\cup Y)=\varnothing$ in the second line. Therefore
\begin{equation}\label{eq:g-diff}
 g(Y\cup Q_v)-g(Y)=\abs\Lambda-\abs{Q_v\cap F(Y)}.
\end{equation}

Now put $d_1=\abs{N_1^+(v)}$ and $d_2=\abs{N_2^+(v)}$, the degrees \emph{in $D$}. By definition of
$Q_v$ the first out-neighborhood loses exactly $Q_v$:
\begin{equation}\label{eq:first-loss}
 \abs{N_{1,D_Y}^+(v)}=d_1-\abs{Q_v}.
\end{equation}
By \eqref{eq:Lambda-sub} and \eqref{eq:confined}, the second out-neighborhood can only lose points
of $\Lambda$ and can only gain points of $Q_v\cap F(Y)$:
\begin{equation}\label{eq:second-change}
 \abs{N_{2,D_Y}^+(v)}\le d_2-\abs\Lambda+\abs{Q_v\cap F(Y)}.
\end{equation}
Substituting \eqref{eq:first-loss} and \eqref{eq:second-change} into \eqref{eq:DY-inequality},
\begin{equation}\label{eq:increment-bound}
 \abs\Lambda-\abs{Q_v\cap F(Y)}
 \overset{\eqref{eq:g-diff}}{=}g(Y\cup Q_v)-g(Y)
 \le \abs{Q_v}+d_2-d_1<\abs{Q_v},
\end{equation}
where the last inequality is the counterexample hypothesis $d_2<d_1$.

\steplabel{Step 5: iterate and sum.}
Start with $Y_0=\varnothing$ and, as long as $Y_i\ne E$, apply Step 3--4 to $Y_i$ to obtain
$Q_{v_i}\ne\varnothing$ with $Q_{v_i}\subseteq E\setminus Y_i$, and set $Y_{i+1}=Y_i\cup Q_{v_i}$.
The sets $Y_i$ strictly increase inside the finite set $E$, so the process stops, and it can only
stop at $Y_r=E$. Each step is performed afresh on the \emph{original} $D$, so minimality applies at
every step. Summing \eqref{eq:increment-bound} over the steps and telescoping,
\[
 g(E)-g(\varnothing)=\sum_{i=0}^{r-1}\bigl(g(Y_{i+1})-g(Y_i)\bigr)<\sum_{i=0}^{r-1}\abs{Q_{v_i}}
 =\abs{E},
\]
because the $Q_{v_i}$ partition $E$. As $g(\varnothing)=0$ and, by \eqref{eq:deficit-via-g},
$g(E)=\abs{F^2(S)\setminus F(S)}$ while $\abs E=\abs{F(S)\setminus S}$, this is
\eqref{eq:deficit}. As $S$ was arbitrary, the lemma is proved.
\end{proof}

\begin{figure}[t]
\centering
\begin{tikzpicture}[
  dot/.style={circle,draw=black!60,fill=#1,inner sep=0pt,minimum size=4mm},
  dot/.default=black!10,
  grp/.style={draw=#1,fill=#1!10,rounded corners=3pt},
  grp/.default=black!40,
  arc/.style={-{Stealth[length=1.6mm]},draw=black!75},
  del/.style={-{Stealth[length=1.6mm]},draw=black!30,dashed},
  lab/.style={font=\scriptsize,text=black!80},
  annot/.style={font=\scriptsize,text=black!65,align=left}
]
% v
\node[circle,draw=black!75,fill=black!10,inner sep=0pt,minimum size=7mm,font=\small] (v) at (0.7,2.8) {$v$};
% three groups of out-neighbours (drawn as boxes, labels above)
\draw[grp=blue!55!black] (2.25,4.05) rectangle (5.05,4.95);
\node[lab,anchor=south] at (3.65,4.99) {$F(v)\cap I$};
\node[dot=blue!25] at (2.85,4.5) {}; \node[dot=blue!25] at (3.65,4.5) {}; \node[dot=blue!25] at (4.45,4.5) {};
\draw[grp=orange!60!black] (2.25,2.35) rectangle (5.05,3.25);
\node[lab,anchor=south] at (3.65,3.29) {$F(v)\cap Y$};
\node[dot=orange!30] at (2.85,2.8) {}; \node[dot=orange!30] at (3.65,2.8) {}; \node[dot=orange!30] at (4.45,2.8) {};
\draw[grp=red!60!black] (2.25,0.65) rectangle (5.05,1.55);
\node[lab,anchor=south] at (3.65,1.59) {$Q_v\subseteq E\setminus Y$};
\node[dot=red!30] at (2.85,1.1) {}; \node[dot=red!30] at (3.65,1.1) {}; \node[dot=red!30] at (4.45,1.1) {};
% kept and deleted arcs out of v
\draw[arc] (v.east) to[out=25,in=180] (2.25,4.45);
\draw[arc] (v.east) to[out=0,in=180]  (2.25,2.8);
\draw[del] (v.south east) to[out=-25,in=180] (2.25,1.15);
\node[font=\large,text=red!55!black] at (1.62,1.72) {$\times$};
% F(Y) and Lambda
\draw[grp=green!45!black] (6.3,2.4) rectangle (9.0,4.4);
\node[lab,anchor=south] at (7.65,4.44) {$F(Y)$};
\node[dot=green!28] at (6.85,3.8) {}; \node[dot=green!28] at (7.65,3.8) {}; \node[dot=green!28] at (8.45,3.8) {};
\node[dot=green!28] at (6.85,2.95) {}; \node[dot=green!28] at (7.65,2.95) {}; \node[dot=green!28] at (8.45,2.95) {};
\draw[grp=red!60!black] (6.55,0.65) rectangle (8.55,1.55);
\node[lab,anchor=south] at (7.55,1.59) {$\Lambda$};
\node[dot=red!35] at (7.1,1.1) {}; \node[dot=red!35] at (8.0,1.1) {};
% arcs into F(Y) and Lambda
\draw[arc] (5.05,2.95) -- (6.25,3.35);
\draw[arc] (5.05,2.65) -- (6.25,3.05);
\draw[arc] (5.05,1.25) -- (6.5,1.25);
\draw[arc] (5.05,0.95) -- (6.5,0.95);
% side annotations
\node[annot,anchor=west] at (9.15,3.55) {only possible new\\ points: $Q_v\cap F(Y)$};
\node[annot,anchor=west] at (8.7,1.1) {counted before,\\ not after the deletion};
% legend, clear of everything
\node[annot,anchor=west] at (0.1,0.52) {dashed: all arcs $S\to E\setminus Y$ are deleted};
\node[annot,anchor=west] at (0.1,0.12) {after the deletion: $F_{D_Y}^2(v)\subseteq I\cup Y\cup F(Y)$};
\end{tikzpicture}

\caption{The deletion step of \cref{lem:deficit}. \emph{Look at} the crossed-out arcs on the left:
 all arcs from $S$ to $E\setminus Y$ are deleted, so the out-neighbors of $v$ that lie in $Q_v$
 disappear, and with them every two-step path that used one of them. \emph{Why it matters:} the
 second out-neighborhood is confined to $I\cup Y\cup F(Y)$ after the deletion
 \eqref{eq:confined}; the points of $\Lambda$ (bottom right) were counted before and are not
 counted after, while the only points that can be gained are the points of $Q_v\cap F(Y)$ ---
 which are also first out-neighbors of $v$, hence are paid for by the first loss.}
\label{fig:deletion}
\end{figure}

\begin{remark}[The ledger, and why the right-hand side is $\abs{F(S)\setminus S}$]\label{rem:ledger}
The four facts of Step 4 are best remembered as a table of what happens to the two neighborhoods of
$v$ when the arcs from $S$ to $E\setminus Y$ are deleted.

\begin{center}
\small
\renewcommand{\arraystretch}{1.25}
\begin{tabular}{@{}llll@{}}
\toprule
quantity & in $D$ & in $D_Y$ & reason \\
\midrule
$\abs{N_1^+(v)}$ & $d_1$ & $d_1-\abs{Q_v}$ & loses exactly $Q_v$, \eqref{eq:first-loss}\\
$\abs{N_2^+(v)}$ & $d_2$ & $\le d_2-\abs\Lambda+\abs{Q_v\cap F(Y)}$ & \eqref{eq:Lambda-sub}, \eqref{eq:confined}, \eqref{eq:new}\\
$g(Y)$ & $g(Y)$ & $g(Y)+\abs\Lambda-\abs{Q_v\cap F(Y)}$ & exact, \eqref{eq:g-diff}\\
\bottomrule
\end{tabular}
\end{center}

Two remarks explain the shape of the lemma.

\emph{(i) The gain is charged to the first out-neighbors.} The new second out-neighbors lie in
$Q_v\cap F(Y)\subseteq Q_v$, and $Q_v$ is precisely the set of first out-neighbors that the
deletion removed. This is why the increment is bounded by $\abs{Q_v}$ and not by $\abs{N_1^+(v)}$:
the deletion pays for the gain out of the first neighborhood, and the deficit $d_2<d_1$ supplies the
strictness.

\emph{(ii) Why $I$ does not appear on the right-hand side.} If the right-hand side of
\eqref{eq:deficit} were $\abs{F(S)}$, the display \eqref{eq:g-diff} would have to account for
$I$ as well, and $I$ is not controlled by the deletion at all. Working with
$F(S)\setminus S=E$ means that the internal part of the image is simply not part of the ledger.
\end{remark}

\subsection{The blowup}\label{sec:blowup}

\begin{ideabox}
\Cref{eq:deficit} has an integer left-hand side and an integer right-hand side, so it gives a
margin of at least $1$. The final argument, however, needs the inequality $\abs{U}+\abs{F^2(U)}<
2\abs{F(U)}$ to survive a comparison in which a term of size at most $n$ (the number of base
vertices) competes with the margin. Replacing every vertex by many twins multiplies the margin by
the block size $m$ while the rounding error stays at $n$: choose $m>n$ and the margin wins.
\end{ideabox}

\begin{lemma}[Transitive-tournament blowup]\label{lem:blowup}
Let $D$ be an oriented graph on $n\ge1$ vertices, all of positive indegree, satisfying
\eqref{eq:deficit} for every nonempty proper subset of $V(D)$. Let $m>n$ be an integer. Replace
every vertex of $D$ by a transitive tournament on $m$ vertices; for every arc $v\to w$ of $D$ put
\emph{all} arcs from the block of $v$ to the block of $w$, and put no other arcs between distinct
blocks. Then the resulting graph $G$ is oriented, all its indegrees are positive, and
\[
 \abs{U}+\abs{F_G^2(U)}<2\abs{F_G(U)}
 \qquad\text{for every }\varnothing\ne U\subsetneq V(G).
\]
\end{lemma}

\inwords{Putting a large transitive tournament in place of each vertex multiplies the subset deficit
by the block size, which beats the rounding error.}

\usedin{Combined with \cref{lem:deficit} it gives \cref{prop:reduction}, hence the hypothesis of
\cref{prop:no-concavity}}

\begin{planbox}
\begin{enumerate}[leftmargin=1.3em,itemsep=1pt,topsep=1pt]
\item \textbf{Internal accounting} (Step~1): inside one block, a nonempty set $T$ of
  labels has image $\tau(T)=\{a+1,\dots,m\}$ where $a=\min T$; hence $\abs T-\abs{\tau(T)}\le1$,
  $\tau^2(T)\subseteq\tau(T)$ and $\abs T-2\abs{\tau(T)}+\abs{\tau^2(T)}\le1$.
\item \textbf{External accounting} (Step~3): if $U$ meets exactly the blocks in $S$, then
  every block indexed by $F(S)$ is \emph{completely} contained in $F_G(U)$, and every block
  indexed by $F^2(S)$ is completely contained in $F_G^2(U)$.
\item \textbf{The ledger} (Step~4): combine the two accounts; the coefficient of $m$ is
  exactly $\abs{F^2(S)\setminus F(S)}-\abs{F(S)\setminus S}\le-1$, and everything else is at most
  $n$.
\item \textbf{The case $S=V(D)$} (Step~5), where the hypothesis \eqref{eq:deficit} is not
  available and a separate argument is needed.
\end{enumerate}
\end{planbox}

\begin{proof}
\steplabel{Step 1: inside a block.}
Write $V(G)=V(D)\times\{1,\dots,m\}$, with the block of $v$ being $\{v\}\times\{1,\dots,m\}$, and
put an arc from $(v,a)$ to $(v,b)$ exactly when $a<b$. Let $\tau$ be the image operator of this
transitive tournament: for $T\subseteq\{1,\dots,m\}$,
\begin{equation}\label{eq:tau}
 \tau(T)=\{b:\ \exists a\in T,\ a<b\}
 =\begin{cases}\{a+1,\dots,m\},&a=\min T,\ T\ne\varnothing,\\ \varnothing,&T=\varnothing.\end{cases}
\end{equation}
For nonempty $T$ with $a=\min T$ we have $T\subseteq\{a,\dots,m\}$ and $\tau(T)=\{a+1,\dots,m\}$,
so
\begin{equation}\label{eq:tau-bounds}
 \abs T-\abs{\tau(T)}\le1,\qquad
 \tau^2(T)\subseteq\tau(T),\qquad
 \abs T-2\abs{\tau(T)}+\abs{\tau^2(T)}\le1 .
\end{equation}
(The inclusion is immediate from \eqref{eq:tau}. The first inequality says that $T$ contains at
most one label below $\tau(T)$; the second that $\tau(T)$ is an interval ending at $m$, so its own
image is smaller still.)

\steplabel{Step 2: the construction is oriented with positive indegrees.}
Inside a block the arcs are transitive, so no loops and no opposite pairs; between two blocks the
arcs all point one way, determined by the arc of $D$; and there is no arc between non-adjacent
blocks. So $G$ is oriented. A vertex $(v,a)$ receives arcs from every block $\{u\}\times\{1,\dots,m\}$
with $u\to v$ in $D$, and every vertex of $D$ has an in-neighbor by hypothesis; so every indegree in
$G$ is positive.

\steplabel{Step 3: which blocks are completely covered.}
Fix a nonempty proper $U\subseteq V(G)$ and let
\begin{equation}\label{eq:support}
 \begin{gathered}
  T_v=\{a:(v,a)\in U\}\subseteq\{1,\dots,m\},\qquad
  S=\{v\in V(D):T_v\ne\varnothing\},\\
  W=F(S)\cup F^2(S).
 \end{gathered}
\end{equation}
\Cref{fig:blowup} shows the three kinds of blocks. First, \emph{every block indexed by $F(S)$ is
completely contained in $F_G(U)$}: if $w\in F(S)$ there is $v\in S$ with $v\to w$, and then
$(v,a)\to(w,b)$ for \emph{all} $a,b$, so the whole block of $w$ is reached. In a block $w\notin
F(S)$ there are no arcs from $U$ into it, so its contribution to $F_G(U)$ is exactly the internal
image $\tau(T_w)$. Therefore
\begin{equation}\label{eq:FG}
 \abs{F_G(U)}=m\,\abs{F(S)}+\sum_{v\in S\setminus F(S)}\abs{\tau(T_v)} .
\end{equation}
Second, \emph{every block indexed by $F^2(S)$ is completely contained in $F_G^2(U)$}: a two-step
path from $U$ whose two steps are external ends anywhere in a block of $F(F(S))=F^2(S)$. A
two-step path with exactly one external step lands in a block of $F(S)$ (if the external step is
the second) or stays inside a block of $F(S)$ (if it is the first, and the block is already full).
A two-step path with no external step contributes $\tau^2(T_v)$ inside the block of $v$. Since
blocks indexed by $F(S)\subseteq W$ are already fully counted,
\begin{equation}\label{eq:FG2}
 \abs{F_G^2(U)}\le m\,\abs W+\sum_{v\in S\setminus W}\abs{\tau^2(T_v)} .
\end{equation}

\steplabel{Step 4: the ledger.}
Subtracting twice \eqref{eq:FG} from \eqref{eq:FG2} and adding $\abs U=\sum_{v\in S}\abs{T_v}$,
\begin{align}
 \abs U-2\abs{F_G(U)}+\abs{F_G^2(U)}
 &\le m\bigl(\abs W-2\abs{F(S)}\bigr)+\sum_{v\in S}\abs{T_v}\notag\\
 &\qquad{}-2\!\!\sum_{v\in S\setminus F(S)}\!\!\abs{\tau(T_v)}
       +\sum_{v\in S\setminus W}\abs{\tau^2(T_v)} .\label{eq:blowup-sharp}
\end{align}
Since $F(S)\subseteq W$, extending the last sum from $S\setminus W$ to $S\setminus F(S)$ only adds
nonnegative terms; this may count an internal second image twice, which is why the display is an
inequality. Now use $\abs{T_v}\le m$ for $v\in S\cap F(S)$ and \eqref{eq:tau-bounds} for
$v\in S\setminus F(S)$:
\begin{align}
 \abs U-2\abs{F_G(U)}+\abs{F_G^2(U)}
 &\le m\bigl(\abs W-2\abs{F(S)}+\abs{S\cap F(S)}\bigr)+\abs{S\setminus F(S)}\notag\\
 &=m\bigl(\abs{F^2(S)\setminus F(S)}-\abs{F(S)\setminus S}\bigr)+\abs{S\setminus F(S)}
 \label{eq:blowup-main}
\end{align}
because $W=F(S)\cup F^2(S)$ gives $\abs W=\abs{F(S)}+\abs{F^2(S)\setminus F(S)}$, and
$\abs{F(S)}-\abs{S\cap F(S)}=\abs{F(S)\setminus S}$.

If $S\ne V(D)$ then $S$ is a nonempty proper subset of $V(D)$, so by hypothesis the parenthesis in
\eqref{eq:blowup-main} is an integer $\le-1$. As $\abs{S\setminus F(S)}\le\abs S\le n<m$,
\[
 \abs U-2\abs{F_G(U)}+\abs{F_G^2(U)}\le -m+\abs{S\setminus F(S)}\le-m+n<0,
\]
which is \eqref{eq:strict-growth} for this $U$.

\steplabel{Step 5: the case $S=V(D)$.}
Here \eqref{eq:deficit} is not available, since $S$ is not proper. But $F(S)=V(D)$: every vertex of
$D$ has an in-neighbor, so every vertex of $D$ is in $F(V(D))$. Hence $F_G(U)=V(G)$, every block
being full. And then $F_G^2(U)=V(G)$ too, because every vertex of $G$ has positive indegree (Step
2). Therefore
\[
 \abs U-2\abs{F_G(U)}+\abs{F_G^2(U)}=\abs U-2\abs{V(G)}+\abs{V(G)}=\abs U-\abs{V(G)}<0,
\]
as $U$ is proper. This is again \eqref{eq:strict-growth}, and the lemma is proved.
\end{proof}

\begin{figure}[t]
\centering
\begin{tikzpicture}[
  base/.style={circle,draw=black!70,fill=black!8,inner sep=0pt,minimum size=5.2mm,font=\scriptsize},
  labd/.style={circle,draw=black!55,fill=black!6,inner sep=0pt,minimum size=3mm},
  arc/.style={-{Stealth[length=1.4mm]},draw=black!70},
  thin/.style={draw=black!30,line width=0.25pt},
  lab/.style={font=\scriptsize,text=black!75},
  pan/.style={font=\scriptsize\bfseries,text=black!55}
]
% ---------------- panel 1: the base graph
\node[base] (u) at (0.7,2.05) {$u$};
\node[base] (v) at (0.7,1.25) {$v$};
\node[base] (w) at (1.75,1.25) {$w$};
\draw[arc] (u) -- (v);
\draw[arc] (v) -- (w);
\node[lab,anchor=north] at (1.05,0.35) {base graph $D$};
% ---------------- panel 2: the blowup of one arc
\draw[rounded corners=3pt,draw=black!35] (3.35,0.55) rectangle (4.25,2.35);
\draw[rounded corners=3pt,draw=black!35] (5.55,0.55) rectangle (6.45,2.35);
\foreach \i in {0,1,2,3}{
  \node[labd] (a\i) at (3.8,0.85+\i*0.42) {};
  \node[labd] (b\i) at (6.0,0.85+\i*0.42) {};
}
\draw[arc] (a0) -- (a3); \draw[arc] (a1) -- (a2);
\draw[arc] (b0) -- (b3); \draw[arc] (b1) -- (b2);
\foreach \i in {0,1,2,3}{ \foreach \j in {0,1,2,3}{ \draw[thin] (a\i) -- (b\j); }}
\node[lab,anchor=north] at (3.8,0.45) {block $v$};
\node[lab,anchor=north] at (6.0,0.45) {block $w$};
\node[lab,anchor=south] at (4.9,2.45) {all arcs $v\to w$};
% ---------------- panel 3: the three kinds of blocks
\draw[fill=black!55,draw=black!55] (7.75,0.55) rectangle (8.35,2.35);
\draw[fill=black!28,draw=black!28] (9.25,0.55) rectangle (9.85,2.35);
\draw[fill=black!8,draw=black!45] (10.75,0.55) rectangle (11.35,2.35);
\node[lab,anchor=north] at (8.05,0.45) {$F(S)$};
\node[lab,anchor=north] at (9.55,0.45) {$F^2(S)\setminus F(S)$};
\node[lab,anchor=north] at (11.05,0.45) {else};
\node[lab,anchor=south] at (9.55,2.45) {the three kinds of blocks};

\end{tikzpicture}

\caption{The blowup $G=D[T_m]$. \emph{Left:} the base graph $D$; a vertex $v$ and an arc $v\to w$.
 \emph{Middle:} each vertex is replaced by a transitive tournament on $m$ labels $1<\dots<m$; the arc
 $v\to w$ becomes a complete set of arcs from block $v$ to block $w$. \emph{Right:} for a set $U$
 meeting the blocks in $S$, the blocks of $F(S)$ are \emph{full} in $F_G(U)$ (dark) and the blocks of
 $F^2(S)\setminus F(S)$ are full in $F_G^2(U)$ (medium); in the remaining blocks only the internal
 image $\tau(T_v)$ is seen. \emph{Why it matters:} the full blocks contribute a multiple of $m$,
 which is where the amplification comes from; the partial blocks contribute at most $n$, which is
 the rounding error.}
\label{fig:blowup}
\end{figure}

\begin{remark}[Reading \eqref{eq:blowup-main}]\label{rem:blowup-read}
The single line \eqref{eq:blowup-main} is the whole point of the blowup, and it is worth seeing why
it is exactly what \cref{eq:deficit} says. The parenthesis
$\abs W-2\abs{F(S)}+\abs{S\cap F(S)}$ counts the blocks as follows: each block of
$F^2(S)\setminus F(S)$ contributes $+1$ (it is full in the second image but not in the first), each
block of $F(S)\setminus S$ contributes $-1$ (it is full in the first image and not in $S$), blocks
of $S\cap F(S)$ contribute $0$ (full in the first image and already in $U$), and all other blocks
contribute $0$ (they are in neither full set). So the parenthesis is precisely
$\abs{F^2(S)\setminus F(S)}-\abs{F(S)\setminus S}$: the deficit of \cref{lem:deficit}, block by
block. The factor $m$ is the number of vertices in a full block, and $\abs{S\setminus F(S)}\le n$
is the price of the blocks that are only partly seen.
\end{remark}

\begin{checkpoint}
Two sentences summarise the section. \emph{Minimality plus arc deletion} produces, for every proper
nonempty $S$, a strict deficit $\abs{F^2(S)\setminus F(S)}<\abs{F(S)\setminus S}$
(\cref{lem:deficit}); the deletion pays for the second-image gain out of the first-image loss, which
is why the deficit is strict. \emph{Blowing up} multiplies that deficit by the block size $m$ and
adds an error of at most $n$, so $m>n$ makes the sum strictly negative
(\cref{lem:blowup}). The next section builds the tool that will convert such inequalities into a
contradiction.
\end{checkpoint}
